\section{GlyphChain distributed-network qualification foundation}
\label{sec:glyphchain-distributed-network-qualification}

This stage adds a repeatable, content-free pre-production exercise for the
GlyphChain transport, partition-recovery and protocol-upgrade controls.  Its
purpose is to make the later physical-device and public-relay trial precise and
reproducible.  A successful local report is evidence that the control paths
execute as designed; it is not evidence that a production network has already
been deployed.

\subsection{Qualification scope and assurance boundary}

The exercise creates four isolated logical profiles representing a home node,
an office node, a virtual-machine node and an intermittently connected node.
Each profile has its own durable state directory.  The exercise covers relay
restart, delayed redelivery, recipient acknowledgement, partition conflict,
explicit recovery, quorum readiness, activation-height enforcement and restart
after a protocol upgrade.

The generated report always sets \texttt{production\_ready=false}.  In
particular, one computer running four isolated profiles does not establish any
of the following:

\begin{itemize}
  \item operation across four independent physical or virtual hosts;
  \item traversal of a residential, carrier-grade or enterprise NAT;
  \item behaviour under a real internet partition or relay-provider outage;
  \item availability of an independently operated public relay;
  \item approval of the malicious-validator response procedure; or
  \item production decentralisation, capacity or latency.
\end{itemize}

These claims remain release gates.  The same exercise and evidence format must
be used again in the representative multi-device environment.

\subsection{Outbound relay and intermittent-node recovery}

The relay exercise uses the bounded \texttt{NATRelayStore}.  An office profile
places a signed-protocol envelope into a queue addressed to an intermittent
profile.  The relay process is reconstructed before the recipient returns,
proving that the queue survives restart.  The recipient pulls the message twice
before acknowledging it.  This models an interrupted client that receives a
message but loses connectivity before acknowledgement.

The relay identifier is content-derived:

\begin{equation}
  r = \operatorname{SHA256}\!\left(
    \operatorname{canon}(c, n_s, n_t, e)
  \right),
\end{equation}

where $c$ is the chain identifier, $n_s$ and $n_t$ are the sender and target
node identifiers, and $e$ is the original protocol envelope.  Repeating the
same put operation produces the same $r$ and does not add a second queue item.
The queue item is removed only after an acknowledgement naming $r$ is received
from the target node.

The exercise checks that the relay envelope contains only bounded network
protocol metadata.  Document bodies, signer email addresses, Legal matter
records and delivery attachments are not relay payloads.  The public relay is
therefore a transport convenience and never a Legal document repository,
validator, finality authority or trust anchor.

\subsection{Partition detection, safety halt and recovery}

Two already validated finality views are introduced for the same height with
different block identifiers.  The classification must be
\texttt{SAFETY\_HALT}.  The durable partition journal binds the transition,
local height, target height and previous event hash.  Restarting the store does
not clear the halt.

Recovery requires all of the following:

\begin{enumerate}
  \item the exact halt identifier;
  \item an identified administrator;
  \item a retained reason for the clearance; and
  \item the exact confirmation phrase
  \texttt{CLEAR GLYPHCHAIN PARTITION SAFETY HALT}.
\end{enumerate}

Clearance moves the node to \texttt{REJOIN\_ASSESSING}, not directly to normal
consensus.  A higher trusted finality view then moves it through
\texttt{BLOCK\_CATCHUP}.  It returns to \texttt{HEALTHY} only after the exact
missing blocks have been materialised and the trusted finality tip has been
reached.  The local qualification report retains only a one-way commitment to
the halt identifier and the number of journal events.

\subsection{Manifest-bound multi-node protocol activation}

Every logical profile independently schedules the same upgrade plan.  The plan
binds the source version, target version, activation height, exact release
manifest digest, validator checkpoint identifier and governance approval
identifier:

\begin{equation}
  p = \operatorname{SHA256}\!\left(
    \operatorname{canon}(c,v_0,v_1,h_a,m,q,g)
  \right).
\end{equation}

Here $m$ is the SHA-256 digest of the exact signed release manifest, $q$ is the
validator-checkpoint identifier and $g$ is the governance-approval identifier.
The exercise fails unless all four profiles derive the same $p$.

Before height $h_a$, each validator advertises its active and supported protocol
range plus $p$ through its signed peer identity exchange.  At activation, the
ready voting power must satisfy

\begin{equation}
  P_{\mathrm{ready}} \ge
  \left\lceil\frac{2P_{\mathrm{total}}}{3}\right\rceil.
\end{equation}

Consensus fails closed as \texttt{UPGRADE\_HALT} if the local release cannot
execute the target protocol or the matching ready power is below quorum.  The
bounded exercise reaches matching quorum, activates version two at the exact
height, completes the plan on each node, reconstructs all four stores and
verifies that the target version remains durable after restart.

\subsection{Evidence report}

The command-line runner hashes the exact release-manifest file rather than
accepting a descriptive release name.  It writes an owner-only JSON report with
exclusive-create semantics.  Existing reports cannot be overwritten, and
symbolic-link destinations are rejected.  The report contains:

\begin{itemize}
  \item the four logical profile names and one-way node commitments;
  \item booleans for each relay, partition and upgrade assertion;
  \item a commitment to the relay and halt identifiers;
  \item the exact release-manifest digest and activation height;
  \item commitments to the checkpoint and governance identifiers;
  \item the explicit local-simulation claim boundary;
  \item the five external production gates; and
  \item a SHA-256 digest over the canonical report body.
\end{itemize}

The report contains no private Legal document, signer identity, message body,
attachment or Legal matter.  Any change to a check, boundary or identifier
invalidates the report digest.

\subsection{Implemented components and verification}

The implementation consists of the reusable qualification engine
\texttt{backend/modules/consensus/network\_qualification.py}, the operator
runner \texttt{backend/scripts/run\_glyphchain\_network\_qualification.py},
the focused regression
\texttt{backend/tests/glyphchain/test\_network\_qualification.py}, and the
content-free evidence file
\texttt{docs/aion/GLYPHCHAIN\_NETWORK\_QUALIFICATION\_REPORT.json}.

The focused network set contains twenty-four passing tests.  Six directly test
the qualification engine and report controls; the remaining eighteen retest the
underlying NAT relay, partition recovery and protocol-upgrade implementations.
The generated evidence reports
\texttt{LOCAL\_PREPRODUCTION\_EXERCISE\_PASSED}, while deliberately retaining
\texttt{production\_ready=false}.

\subsection{Remaining production qualification}

Before distributed-network approval, Tessaris must still run this procedure on
independent representative home, office, VM and intermittent devices; exercise
real NAT, reconnection, delay, relay outage and network partition conditions;
activate an exact signed release across those devices; validate an
independently operated production-like public relay; and obtain operational and
security approval for malicious-validator and emergency-removal handling.  A
local report must never be relabelled as live deployment evidence.

